% Part: first-order-logic
% Chapter: completeness
% Section: lindenbaums-lemma

\documentclass[../../../include/open-logic-section]{subfiles}

\begin{document}

\iftag{FOL}
      {\olfileid[id]{fol}{com}{lin}}
      {\olfileid[id]{pl}{com}{lin}}

\olsection{Lema Lindenbaum}

\begin{explain}
Sekarang kita membuktikan suatu lema yang menunjukkan bahwa setiap
himpunan konsisten yang terdiri atas !!{sentence}s termuat dalam suatu
himpunan kalimat yang tidak hanya konsisten, tetapi juga !!{complete}.
Dalam pembuktiannya, satu !!{sentence} ditambahkan pada setiap
langkah sambil dijamin bahwa himpunan tersebut tetap konsisten. Kita
melakukannya sedemikian rupa sehingga untuk setiap $!A$, $!A$ atau
$\lnot !A$ ditambahkan pada suatu tahap. Gabungan semua tahap dalam
konstruksi itu kemudian memuat $!A$ atau negasinya~$\lnot !A$, sehingga
bersifat lengkap. Gabungan tersebut juga konsisten karena pada setiap
tahap kita memastikan agar tidak menimbulkan inkonsistensi.
\end{explain}

\begin{lem}[Lema Lindenbaum]
\ollabel{lem:lindenbaum} Setiap himpunan konsisten~$\Gamma$ dalam
bahasa~$\Lang{L}$ dapat diperluas menjadi suatu himpunan~$\Gamma^*$ yang
!!a{complete} dan konsisten.
\end{lem}

\begin{proof}
Misalkan $\Gamma$ konsisten. Misalkan $!A_0$, $!A_1$,
\dots{} merupakan suatu enumerasi semua !!{sentence}s dalam~$\Lang L$.
Definisikan $\Gamma_0 = \Gamma$, dan
\[
\Gamma_{n+1} =
\begin{cases}
\Gamma_n \cup \{ !A_n \} & \textrm{jika $\Gamma_n \cup \{!A_n\}$
  konsisten;} \\
\Gamma_n \cup \{ \lnot !A_n \} & \textrm{jika tidak.}
\end{cases}
\]
Misalkan $\Gamma^* = \bigcup_{n \geq 0} \Gamma_n$.

Setiap $\Gamma_n$ konsisten: $\Gamma_0$ konsisten menurut hipotesis dan
definisi~$\Gamma_0$.
Jika $\Gamma_{n+1} = \Gamma_n \cup \{!A_n\}$, hal ini karena himpunan
yang disebut terakhir itu konsisten. Jika tidak, $\Gamma_{n+1} =
\Gamma_n \cup \{\lnot !A_n\}$. Kita harus memastikan bahwa
$\Gamma_n \cup \{\lnot !A_n\}$ konsisten. Andaikan himpunan itu tak
konsisten. Maka \emph{kedua} himpunan $\Gamma_n \cup \{!A_n\}$ dan
$\Gamma_n \cup \{\lnot !A_n\}$ tak konsisten. Hal ini berarti bahwa
$\Gamma_n$ tak konsisten menurut
\iftag{FOL}{%
  \tagrefs{prfAX/{fol:axd:prv:prop:provability-exhaustive},
    prfSC/{fol:seq:prv:prop:provability-exhaustive},
    prfND/{fol:ntd:prv:prop:provability-exhaustive},
    prfTab/{fol:tab:prv:prop:provability-exhaustive}}}{%
  \tagrefs{prfAX/{pl:axd:prv:prop:provability-exhaustive},
    prfSC/{pl:seq:prv:prop:provability-exhaustive},
    prfND/{pl:ntd:prv:prop:provability-exhaustive},
    prfTab/{pl:tab:prv:prop:provability-exhaustive}}},
bertentangan dengan hipotesis induksi.

Untuk setiap~$n$ dan setiap $i < n$, $\Gamma_i \subseteq \Gamma_n$.
Hal ini mengikuti melalui induksi sederhana pada~$n$. Untuk $n=0$,
tidak terdapat $i < 0$, sehingga klaim tersebut berlaku secara otomatis.
Untuk langkah induksi, andaikan klaim itu benar bagi~$n$. Kita tunjukkan
bahwa jika $i < n+1$, maka $\Gamma_i \subseteq \Gamma_{n+1}$. Menurut
konstruksi, $\Gamma_{n+1} = \Gamma_n \cup \{!A_n\}$ atau
$= \Gamma_n \cup \{\lnot !A_n\}$. Jadi, $\Gamma_n \subseteq
\Gamma_{n+1}$. Jika $i < n+1$, maka $\Gamma_i \subseteq \Gamma_n$
menurut hipotesis induksi (jika $i < n$), atau menurut fakta sepele bahwa
$\Gamma_n \subseteq \Gamma_n$ (jika $i = n$). Kita memperoleh
$\Gamma_i \subseteq \Gamma_{n+1}$ melalui transitivitas
relasi~$\subseteq$.

Dari hasil ini, diperoleh bahwa $\Gamma^*$ konsisten. Alasannya sebagai
berikut: misalkan $\Gamma' \subseteq \Gamma^*$ berhingga. Jika
$\Gamma'=\varnothing$, maka $\Gamma'$ konsisten secara sepele. Jika
$\Gamma'\ne\varnothing$, setiap $!B \in \Gamma'$ juga berada
dalam~$\Gamma_i$ untuk suatu~$i$. Misalkan $n$ adalah indeks terbesar di
antara indeks-indeks tersebut. Karena $\Gamma_i \subseteq \Gamma_n$ jika
$i \le n$, setiap $!B \in \Gamma'$ juga $\in \Gamma_n$, yakni $\Gamma'
\subseteq \Gamma_n$, dan $\Gamma_n$ konsisten. Jadi, setiap himpunan
bagian berhingga $\Gamma' \subseteq \Gamma^*$ konsisten.
Menurut \iftag{FOL}{%
  \tagrefs{prfAX/{fol:axd:ptn:prop:proves-compact},
    prfSC/{fol:seq:ptn:prop:proves-compact},
    prfND/{fol:ntd:ptn:prop:proves-compact},
    prfTab/{fol:tab:ptn:prop:proves-compact}}}{%
  \tagrefs{prfAX/{pl:axd:ptn:prop:proves-compact},
    prfSC/{pl:seq:ptn:prop:proves-compact},
    prfND/{pl:ntd:ptn:prop:proves-compact},
    prfTab/{pl:tab:ptn:prop:proves-compact}}}, $\Gamma^*$ konsisten.

Setiap !!{sentence} dalam $\Frm[L]$ muncul dalam daftar yang digunakan
untuk mendefinisikan~$\Gamma^*$. Jika $!A_n \notin \Gamma^*$, penyebabnya
adalah bahwa $\Gamma_n \cup \{!A_n\}$ tak konsisten. Namun, dalam hal
itu $\lnot !A_n \in \Gamma^*$, sehingga $\Gamma^*$ !!{complete}.
\end{proof}

\end{document}
